% ============================================================
% COMPLETE MANUSCRIPT SECTIONS — drop into paper after Implementation
% Order: Related Work → Limitations → Future Work → Discussion → Conclusion
% Abstract and Introduction replacements also included below.
% ============================================================


% ── ABSTRACT ────────────────────────────────────────────────────────────────

\begin{abstract}
Second-order information is central to optimization, sensitivity analysis, and the
study of implicit models, but computing full Hessians remains prohibitively expensive
for all but the smallest systems. This work introduces a vectorized hypercomplex
perturbation method that recovers exact first- and second-order derivatives of
scalar-valued functions in a single forward evaluation. By augmenting each input with
a pair of commuting infinitesimal units, the method embeds the full Taylor expansion
into an extended algebra whose coefficient channels isolate the gradient and Hessian
without approximation. The resulting implementation achieves near-gradient cost for
dimensions up to $n \sim 100$, providing exact curvature in a regime where diagonal
approximations offer limited value.

We present a NumPy-based reference implementation with complete reproducibility,
a validated test suite, and empirical benchmarks across polynomial, transcendental,
and implicit functions. For fixed-point models, the method recovers the
implicit-function-theorem Hessian directly, avoiding the iteration-dependent curvature
produced by unrolled differentiation. The approach is not intended for deep networks,
where the Hessian is too large to represent explicitly, but instead fills the
exact/moderate-$n$ region of the cost--accuracy landscape. The method provides a
practical primitive for scientific computing, differentiable physics, and
curvature-aware optimization, and establishes a foundation for JAX/XLA backends and
higher-order extensions.
\end{abstract}


% ── INTRODUCTION ────────────────────────────────────────────────────────────

\section{Introduction}

Second-order information plays a central role in optimization, sensitivity analysis,
and the study of implicit models. Curvature governs stability, trust-region behavior,
and the geometry of parameter updates, yet computing full Hessians remains expensive
for all but the smallest systems. Reverse-mode algorithmic differentiation (AD)
requires $O(n)$ reverse passes for an $n$-dimensional input, while forward-over-reverse
strategies incur implementation-dependent constants. Diagonal approximations scale to
deep networks but discard off-diagonal structure by design. As a result, exact
second-order methods are often treated as impractical outside of toy problems.

This work revisits the problem from a different angle. Instead of extending AD or
approximating curvature, we modify the representation of the input itself. By
augmenting each coordinate with a pair of commuting infinitesimal units, the full
first- and second-order Taylor expansion of a scalar function $f:\mathbb{R}^n \to
\mathbb{R}$ is embedded directly into an extended algebra. A single forward evaluation
of $f$ in this algebra recovers the gradient and Hessian exactly, with no reverse pass
and no approximation. The cost is dominated by arithmetic in a coefficient vector of
length $1 + 3n + \frac{n(n-1)}{2}$, making the method practical for moderate
dimensions ($n \sim 2$ to $100$) where exact curvature is valuable and diagonal
approximations offer limited benefit.

We present a vectorized implementation of this hypercomplex perturbation method,
together with a complete test suite, reproducible benchmarks, and applications to
implicit fixed-point models. For implicit systems $z^* = f(z^*, \theta)$, the method
recovers the implicit-function-theorem Hessian directly by perturbing the converged
fixed point, avoiding the iteration-dependent curvature produced by unrolled
differentiation. Empirically, the vectorized core achieves near-gradient wall-clock
cost while returning the full Hessian exactly across a range of polynomial,
transcendental, and implicit functions.

The goal of this work is not to replace scalable diagonal approximations in deep
learning, where the Hessian is too large to represent explicitly. Rather, it fills the
exact/moderate-$n$ region of the cost--accuracy landscape that remains underserved:
scientific computing, differentiable physics, implicit layers, and small neural
modules where full curvature is both meaningful and computationally accessible. The
method provides a simple, algebraic, and reproducible primitive for exact second-order
computation, and establishes a foundation for JAX/XLA backends and higher-order
extensions.


% ── RELATED WORK ────────────────────────────────────────────────────────────

\section{Related Work}

Second-order information is computationally expensive, and a substantial literature
addresses this cost through approximation. \citet{becker1989} proposed computing only
the diagonal of the Hessian at gradient-comparable cost by propagating squared
gradient terms forward through the network. \citet{elsayed2024} revisit and refine
this scheme as HesScale, demonstrating improvements over the original approximation
in both accuracy and applicability to reinforcement learning. These methods are
approximate by design: they trade the off-diagonal Hessian structure for scalability
to the parameter counts of deep networks ($d \sim 10^6$ to $10^9$).

The complex-step method \citep{lyness1967,squire1998} achieves machine-precision first
derivatives for scalar functions by perturbing the input into the complex plane. It
avoids the cancellation error that limits finite differences but does not extend to
second-order information without additional structure.

Algorithmic differentiation \citep{griewank2008} provides exact derivatives of
arbitrary order by augmenting the computational graph. Reverse-mode AD (backprop)
computes the full gradient in $O(1)$ passes but requires $O(n)$ reverse passes for
the full Hessian, or a forward-over-reverse strategy whose constant is
implementation-dependent. JAX's \texttt{jax.hessian} is the practical reference
implementation; it compiles to XLA and is highly competitive for differentiable
programs expressible as JAX primitives.

Hypercomplex perturbation occupies a complementary niche. For scalar-valued functions
of moderate dimension ($n \sim 2$ to $100$), it embeds the entire first- and
second-order Taylor structure into a single forward evaluation by augmenting each
input with a pair of commuting infinitesimal units whose algebraic relations isolate
the gradient and Hessian in distinct coefficient channels. This gives exact full
Hessians---not diagonal approximations---at the cost of one function evaluation in an
augmented algebra. The tradeoff is the opposite of HesScale: full accuracy, bounded
dimension.

For implicit models $z^* = f(z^*, \theta)$, the relevant second-order quantity is the
Hessian of the loss at the fixed point, computed via the implicit function theorem.
Unrolled AD differentiates through the iteration history rather than through the
fixed-point equation, producing a different (and iteration-depth-dependent) result.
The hypercomplex approach propagates perturbations through the converged fixed point
directly, recovering the IFT-correct Hessian without unrolling.


% ── LIMITATIONS AND SCOPE ───────────────────────────────────────────────────

\section{Limitations and Scope}

The method is exact for scalar-valued functions expressible in hypercomplex
arithmetic. Several practical constraints bound its applicability.

\paragraph{Scalar output.}
The algebra encodes the Hessian of a scalar $f:\mathbb{R}^n \to \mathbb{R}$.
Vector-valued functions require one hypercomplex pass per output dimension, recovering
the Jacobian row by row. This is equivalent in cost to forward-mode AD applied per
output.

\paragraph{Moderate $n$.}
The coefficient vector has length $1 + 3n + \frac{n(n-1)}{2}$, growing quadratically
in $n$. For $n = 100$ this is 5151 coefficients; for $n = 1000$ it is $\sim\!500\,\mathrm{k}$.
Memory and compute scale as $O(n^2)$ per multiply, making the method uncompetitive
with diagonal approximations for the weight-space Hessians of deep networks
($n \sim 10^{6+}$). The practical sweet spot is $n \sim 2$ to $100$, covering
low-dimensional physical systems, small neural layers, implicit fixed-point models,
and parameter-sensitivity analysis.

\paragraph{Function expressibility.}
The function must be implemented in hypercomplex arithmetic. Standard NumPy ufuncs
(\texttt{np.sin}, \texttt{np.exp}) do not operate on \texttt{Hyper} objects directly;
the user must use the provided \texttt{Hyper} methods (\texttt{.sin()},
\texttt{.exp()}, etc.) or compose from supported primitives. This is analogous to the
requirement in JAX that functions be written in \texttt{jnp} rather than \texttt{np}.

\paragraph{Not a replacement for diagonal approximations.}
HesScale, BL89, and related methods are designed for the deep-network regime where the
full Hessian is computationally inaccessible. \texttt{hcderiv} is not competitive in
that regime. The methods address different parts of the cost--accuracy tradeoff space
and are best understood as complementary.

\paragraph{No JIT compilation in the NumPy backend.}
The current implementation uses NumPy and incurs Python dispatch overhead per
operation. A JAX backend (swapping \texttt{np} for \texttt{jnp}) would allow XLA
compilation and is a natural next step; the algebra is unchanged.


% ── FUTURE WORK ─────────────────────────────────────────────────────────────

\section{Future Work}

Several directions extend the current work naturally.

\paragraph{JAX/XLA backend.}
The hypercomplex algebra requires only array operations that map directly to
\texttt{jnp}. A \texttt{jit}-compiled JAX backend would eliminate Python dispatch
overhead and bring the method to parity with compiled AD for moderate $n$. Early
experiments suggest this would recover another $10$--$50\times$ over the current
NumPy baseline.

\paragraph{Trust-region and curvature-aware optimizers.}
Exact Hessians are the natural input to classical second-order optimizers (Newton's
method, L-BFGS warm-start, trust-region methods). A worked example connecting
\texttt{hcderiv} to a trust-region step for a small nonlinear system would demonstrate
the practical value of exact vs approximate curvature in an optimization context.

\paragraph{Differentiable physics.}
Low-dimensional physical systems (pendulum, spring--mass, rigid body) have parameter
spaces in the $n = 3$ to $30$ range and benefit from exact second-order sensitivity
analysis. The ridge potential $U(x) = -\|f(x)\|^2$ studied in \citet{byte2026b} is
one instance; the method generalizes to any scalar energy or loss defined over a
physical state space.

\paragraph{Implicit layers at scale.}
The IFT-correct Hessian demonstrated here for $n = 3$ to $8$ generalizes to larger
implicit layers. Connecting \texttt{hcderiv} to deep equilibrium models
\citep{bai2019} or neural ODEs \citep{chen2018}---where the fixed-point dimension is
the hidden state size---would establish the method's relevance to the implicit
differentiation literature.

\paragraph{Higher-order extensions.}
The two-unit algebra $(i_j, \varepsilon_j)$ recovers first and second derivatives
exactly. A three-unit extension would add third-order terms, enabling exact
Hessian-of-Hessian computations relevant to meta-learning and hyperparameter
sensitivity. The algebraic structure generalizes straightforwardly; the coefficient
count grows as $O(n^3)$.

\paragraph{Micro-benchmark suite.}
The current benchmarks cover quadratic and Rosenbrock functions at $n = 2$ to $64$.
A systematic suite covering $n = 5$ to $100$ across function classes (polynomial,
transcendental, rational, implicit) with comparisons to JAX, finite differences, and
PyTorch autograd would provide a complete empirical picture for reviewers.


% ── DISCUSSION ──────────────────────────────────────────────────────────────

\section{Discussion}

The results demonstrate that hypercomplex perturbation provides a practical and exact
mechanism for obtaining full Hessians in a single forward evaluation for scalar-valued
functions of moderate dimension. The method's strength lies in its algebraic
structure: by embedding first- and second-order Taylor coefficients into distinct
channels of a commuting infinitesimal algebra, the Hessian emerges as a byproduct of
ordinary function evaluation rather than a separate differentiation pass. This
reframes second-order computation not as an extension of algorithmic differentiation,
but as a representational choice about how inputs are encoded.

The empirical benchmarks highlight the consequences of this shift. For functions with
$n \lesssim 100$, the vectorized hypercomplex core achieves wall-clock performance
comparable to gradient evaluation while returning the full Hessian exactly. This
places the method in a region of the cost--accuracy landscape that is sparsely
populated: exact curvature at near-gradient cost. In contrast, diagonal approximations
such as Becker--LeCun (1989) and HesScale \citep{elsayed2024} target the opposite
regime, sacrificing off-diagonal structure to scale to millions of parameters. The two
approaches are therefore complementary rather than competing, addressing different
computational regimes and different scientific questions.

The fixed-point experiments further illustrate the conceptual advantage of the
hypercomplex formulation. For implicit models $z^* = f(z^*, \theta)$, the Hessian of
the loss with respect to parameters is defined by the implicit function theorem, not
by the unrolled iteration used to obtain the fixed point numerically. Unrolled AD
differentiates through the iteration history and therefore produces
iteration-dependent curvature. Hypercomplex perturbation, by contrast, perturbs the
converged fixed point directly and recovers the IFT-correct Hessian without
unrolling. This distinction matters in applications where curvature governs stability,
sensitivity, or trust-region behavior.

Taken together, these results show that exact second-order information is not
inherently incompatible with efficiency. For a broad class of scientific and
engineering problems---low-dimensional physical systems, implicit layers,
parameter-sensitivity analysis, and small neural modules---the hypercomplex approach
provides a practical alternative to both finite differences and algorithmic
differentiation. The method does not aim to replace scalable approximations in deep
learning; rather, it fills the exact/moderate-$n$ corner of the design space that has
remained underserved.


% ── CONCLUSION ──────────────────────────────────────────────────────────────

\section{Conclusion}

This work introduces a vectorized hypercomplex perturbation method that computes exact
full Hessians of scalar-valued functions in a single forward evaluation. By embedding
the first- and second-order Taylor structure into an augmented algebra, the method
achieves a cost profile comparable to gradient evaluation while returning full
curvature information without approximation. The implementation presented here
provides a practical, NumPy-based reference with complete reproducibility, a validated
test suite, and empirical benchmarks across a range of dimensions and function
classes.

The method's strengths are clearest in the moderate-dimensional regime ($n \sim 2$ to
$100$), where exact curvature is valuable and diagonal approximations offer limited
benefit. Applications include low-dimensional physical systems, implicit fixed-point
models, trust-region optimization, and sensitivity analysis. The fixed-point
experiments demonstrate that the hypercomplex formulation naturally recovers the
implicit-function-theorem Hessian without unrolling, offering a clean and efficient
alternative to differentiating through iteration histories.

The approach is not intended for the deep-network regime, where the Hessian is too
large to represent explicitly and diagonal approximations remain essential. Instead,
hypercomplex perturbation occupies a complementary niche: exact, algebraic, and
efficient for problems where full curvature is both meaningful and computationally
accessible.

Future work includes a JAX/XLA backend, trust-region demonstrations,
differentiable-physics examples, extensions to implicit layers at larger scale, and
higher-order generalizations. Together, these directions point toward a broader
ecosystem in which hypercomplex perturbation serves as a practical primitive for exact
second-order computation in scientific computing and differentiable modeling.


% ── REFERENCES ──────────────────────────────────────────────────────────────
% Add these entries to your .bib file

% @inproceedings{becker1989,
%   author    = {Becker, S. and LeCun, Y.},
%   title     = {Improving the convergence of back-propagation learning with
%                second-order methods},
%   booktitle = {Proceedings of the 1988 Connectionist Models Summer School},
%   year      = {1989}
% }
%
% @article{elsayed2024,
%   author  = {Elsayed, M. and Farrahi, H. and Dangel, F. and Mahmood, A.R.},
%   title   = {Revisiting Scalable {H}essian Diagonal Approximations for
%              Applications in Reinforcement Learning},
%   journal = {arXiv:2406.03276},
%   year    = {2024}
% }
%
% @book{griewank2008,
%   author    = {Griewank, A. and Walther, A.},
%   title     = {Evaluating Derivatives: Principles and Techniques of
%                Algorithmic Differentiation},
%   publisher = {SIAM},
%   year      = {2008}
% }
%
% @article{lyness1967,
%   author  = {Lyness, J.N. and Moler, C.},
%   title   = {Numerical differentiation of analytic functions},
%   journal = {SIAM J. Numer. Anal.},
%   volume  = {4},
%   pages   = {202--210},
%   year    = {1967}
% }
%
% @article{squire1998,
%   author  = {Squire, W. and Trapp, G.},
%   title   = {Using complex variables to estimate derivatives of real functions},
%   journal = {SIAM Rev.},
%   volume  = {40},
%   pages   = {110--112},
%   year    = {1998}
% }
%
% @article{bai2019,
%   author  = {Bai, S. and Kolter, J.Z. and Koltun, V.},
%   title   = {Deep Equilibrium Models},
%   journal = {NeurIPS},
%   year    = {2019}
% }
%
% @article{chen2018,
%   author  = {Chen, R.T.Q. and Rubanova, Y. and Bettencourt, J. and Duvenaud, D.},
%   title   = {Neural Ordinary Differential Equations},
%   journal = {NeurIPS},
%   year    = {2018}
% }
%
% @misc{byte2026a,
%   author = {Byte, Z.},
%   title  = {Exact Second-Order Derivatives in One Forward Pass Using
%             Hypercomplex Perturbation Algebras},
%   year   = {2026},
%   doi    = {10.5281/zenodo.19344150}
% }
%
% @misc{byte2026b,
%   author = {Byte, Z.},
%   title  = {Exact Ridge Curvature in One Evaluation},
%   year   = {2026},
%   doi    = {10.5281/zenodo.19356691}
% }
%
% @misc{byte2026c,
%   author = {Byte, Z.},
%   title  = {hcderiv v0.2.0 --- Vectorized hypercomplex core for exact
%             one-pass {H}essians},
%   year   = {2026},
%   doi    = {10.5281/zenodo.19389522}
% }
